\documentclass[10pt,a4paper]{paper}
\usepackage[T1]{fontenc}
\usepackage[spanish]{babel}
\title{TP1}
\begin{document}

\maketitle
	
	
	
	
		
\paragraph{1.a}
	
\begin{itemize}
	\item Mencione y explique las características más relevantes de GNU/Linux.


	\item GNU/Linux es un sistema operativo libre que surge de la combinación del núcleo Linux con el sistema GNU a principios de los 90. 
	
	Esta compuesto de principalmente por un kernel, un sistema de archivos jerárquico y un interprete de comandos. El kernel es monolítico híbrido. Es decir se comporta como un único programa en memoria, pero tiene capacidad de cargar y descargar módulos mientras se ejecuta. El kernel se distribuye bajo la licencia libre GPLv2.
	
	Es un sistema operativo multiusuario y multitarea escrito principalmente en C. Esto brinda la posibilidad de soportar varias arquitecturas.
	
	El sistema de archivos provee un modelo de permisos basado en tres grupos, el usuario, su grupo y los demás.

\end{itemize}


		
\paragraph{1.b}

\begin{itemize}
	\item Mencione otros sistemas operativos y compárelos con GNU/Linux en cuanto a los puntos mencionados en el inciso a.
	
	\item Windows es un sistema operativo comercial privativo de código cerrado. Además, Windows es desarrollado principalmente por Microsoft, tanto el kernel como sus utilidades de usuario. El kernel de Windows tiene incorporado más elementos que en Linux son programas de usuario. Pero es también híbrido. Windows esta escrito en C principalmente pero al ser de código cerrado su portabilidad es más limitada en la practica. El sistema de archivos que provee Windows es muy diferente al modelo de Linux en estructura y funcionalidad. Windows soporta FAT, exFAT, y NTFS pero Linux tiene una capa de abstracción, VFS, que permite incorporar nuevos y soporta bastantes.
	
	\item Android esta basado en el kernel Linux y es desarrollado por Google y otros, si bien es de código abierto, esta optimizado para las pantallas táctiles y dispositivos ARM. Esto limita su portabilidad en la practica si ademas tenemos en cuenta que sus utilidades de usuario son limitadas. Android usa las utilidades de usuario \texttt{Toybox}, entre otras, que se caracterizan por ser bastante mínimas en comparación a sistemas GNU.
\end{itemize}

	


\paragraph{1.c}

\begin{itemize}
	\item ¿Qué es GNU?
	
	\item GNU es un sistema operativo tipo Unix. desarrollado para el proyecto GNU y bajo el auspicio de la FSF. Formado principalmente de software libre bajo términos \textit{copyleft} como GPL. GNU significa \textit{GNU is not Unix}, que destaca la diferencia entre GNU y Unix, siendo el primero libre y el segundo privativo. El kernel del proyecto, GNU Hurd, no llegó a completarse. Algunos componentes populares son \texttt{Bash, coreutils, gcc y glibc}.
\end{itemize}


\paragraph{1.d}

\begin{itemize}
	\item Indique una breve historia sobre la evolución del proyecto GNU.
	
	\item Richard Stallman inicia el proyecto en 1983. En 1985 crea la FSF que proporcionará soporte técnico y legal al proyecto GNU. Para el 1990 se han creado GCC, Bash e Emacs y la mayor parte de bibliotecas y utilidades que conformarían el SO, pero no se tiene un kernel funcional. Durante los 80 el proyecto GNU evalúa, e intenta adaptar kernels existentes. A principios de los 90 se opta por adaptar Mach de la Universidad Carnegie Mellon, que se denominaría Hurd. Sin embargo, este se demoraría y su adopción es limitada. De forma autónoma empieza el desarrollo de un kernel por el informático finlandés Linus Torvalds en 1991. Torvalds adoptaría la licencia GPL para su proyecto en el año 1992. Esto precipita el desarrollo del kernel que tomaría el nombre Linux, y que da lugar a la adopción de este junto a las herramientas GNU integradas en un SO libre que predomina hasta hoy.
\end{itemize}


\paragraph{1.e}

\begin{itemize}
	\item Explique qué es la multitarea, e indique si GNU/Linux hace uso de ella.
	
	\item La multitarea es la capacidad de un SO de conmutar y administrar múltiples tareas, es decir, procesos e hilos, de forma concurrente. GNU/Linux es multitarea. Específicamente usa multitarea apropiativa: el kernel puede interrumpir una tarea en ejecución para asignar tiempo de CPU a otra.
\end{itemize}

\paragraph{1.f}

\begin{itemize}
	\item ¿Qué es POSIX?
	
	\item POSIX es una norma de la IEEE que define una interfaz estándar del SO y su entorno, incluyendo al intérprete de comandos y programas de utilidades con el objetivo de soportar la portabilidad de aplicaciones, esta inspirada en interfaces, intérpretes de comandos y utilidades comunes en versiones de UNIX. Hoy en día muchos SO cumplen parcial o totalmente con algún estándar POSIX.
\end{itemize}


\paragraph{2.a}

\begin{itemize}
	\item ¿Qué es una distribución de GNU/Linux? Nombre al menos 4 distribuciones de GNU/Linux y cite diferencias básicas entre ellas.
	
	\item Una distribución de GNU/Linux es un gran conjunto de programas binarios empaquetados para ser distribuidos. En una distribución generalmente se incluyen programas de instalación que facilitan el proceso de configuración e instalación del SO. 
	
	Algunas distribuciones son Debian, Ubuntu, Fedora, Arch. las principales diferencias son sus administradores de paquetes, los mecanismos de actualización y las metodologías de desarrollo. 
\end{itemize}


\paragraph{2.b}

\begin{itemize}
	\item ¿En qué se diferencia una distribución de otra?
	
	\item Generalmente usan gestores de paquetes diferentes. Tienen un modelo de lanzamientos diferente, es diferente el propósito que cumplen, el método de instalación, o su filosofía. 
\end{itemize}


\paragraph{2.c}

\begin{itemize}
	\item ¿Qué es Debian? Acceda al sitio https://www.debian.org/ e indique cuáles son los objetivos del proyecto y una breve cronología del mismo.
	
	\item Debian es una distribución GNU/Linux desarrollada por el proyecto Debian. Entre sus objetivos se encuentran: el crear un sistema operativo completamente libre y accesible a todos, mantener el software totalmente libre, ser estable, seguro y de alta calidad, ser universal y servir de base para otras distribuciones y contribuir de vuelta a la comunidad de software libre.
	
	Históricamente Debian, fundado en 1993 por Ian Murdock tiene su primera versión 1.1 llamada Debian Buzz en el año 1996. Actualmente nos encontramos en su versión 12 Debian Bookworm que fue lanzada en 2023. 
\end{itemize}


\paragraph{3.a}

\begin{itemize}
	\item Nombre cuáles son los componentes fundamentales de GNU/Linux.
	
	
	\item Un SO GNU/Linux está compuesto principalmente por el kernel Linux, una biblioteca C estándar como glibc, utilidades de usuario GNU, un intérprete de comandos como Bash, un gestor de arranque como GRUB, un sistema de inicio como systemd, un sistema de archivos,  y otros componentes de bajo nivel como PulseAudio, Mesa, etc. según su necesidad.
\end{itemize}


\paragraph{3.b}

\begin{itemize}
	\item Mencione y explique la estructura básica del Sistema Operativo GNU/Linux.
	
	
	\item Un SO GNU/Linux tiene una estructura en capas. En la capa más interna está el kernel, que se ejecuta en modo privilegiado. Por encima de este están los manejadores, que a veces son parte del kernel. Encima está la biblioteca C estándar. Esta última está en el espacio de usuario junto con las demás aplicaciones de usuario como las utilidades GNU, bash y demás comandos.
\end{itemize}



\paragraph{4.a}

\begin{itemize}
	\item ¿Cuáles son las funciones principales del kernel?
	
	\item Las funciones principales del kernel son:
	\begin{itemize}
		\item Proveer a los usuarios de abstracciones útiles por medio de las llamadas al sistema.
		
		\item Administrar los recursos de hardware de forma equitativa, eficiente y segura según se requiera.
		
		\item Gestionar peticiones de entrada/salida de los procesos.
	\end{itemize}
\end{itemize}


\paragraph{4.b}

\begin{itemize}
	\item ¿Es posible tener más de un kernel de GNU/Linux instalado en la misma máquina? (ayuda /boot).
	
	
	\item Si es posible, como ejemplo en una distribución Ubuntu en el directorio /boot se puede tener: \texttt{vmlinuz, vmlinuz-7.0.0-29-generic, vmlinuz-7.0.0-30-generic, vmlinuz.old}. Todas son imágenes del kernel.
\end{itemize}


\paragraph{4.c}

\begin{itemize}
	\item ¿En qué directorio se encuentra ubicado el kernel?
	
	
	\item Como se insinuó está en \texttt{/boot}. Estas entradas son las que es posible seleccionar desde el menú que presenta normalmente GRUB.
\end{itemize}



\paragraph{5.a}

\begin{itemize}
	\item ¿Qué es y cuáles son las funciones del intérprete de comandos o \textit{shell}?
	
	
	\item El intérprete de comandos es un programa que presenta una interfaz al usuario para poder interactuar con el SO. Una shell permite ejecutar programas en el SO, como comandos y aplicaciones. Además Bash, por ejemplo, soporta redireccionamiento, tuberías, alias, historial, etc.
\end{itemize}

\paragraph{5.b}

\begin{itemize}
	\item Investigue y mencione al menos 3 intérpretes de comandos que existen para GNU/Linux y compárelos entre ellos.
	
	
	\item Tres shells son csh, zsh y bash. Hacer una comparacion relevante de diferencias me es imposible en este momento.:(
\end{itemize}



\paragraph{5.c}

\begin{itemize}
	\item ¿Dónde se ubican (path) los comandos propios y externos al Shell?
	
	\item Los comandos propios de la shell son parte de esta y se ejecutan internamente,  por ej. cd, alias, etc. 
	
	
	En cambio los externos se ubican en archivos que la shell busca en los directorios de la v. de entorno \texttt{PATH}, típicamente en \texttt{/bin}, \texttt{/usr/bin}. Entre otros.
	
	
\end{itemize}


\paragraph{5.d}

\begin{itemize}
	\item ¿Por qué considera que la Shell no es parte del Kernel de GNU/Linux?
	
	\item La shell no forma parte del kernel por razones de diseño. No existe necesidad de que esta forme parte del kernel. 
	
	Para empezar, la shell no interactúa con hardware, puede usar glibc y ejecutar sus comandos, etc. Es decir, la shell es un programa de usuario que aprovecha las abstracciones del kernel y otras bibliotecas, por esto no es necesario que forme parte de este.	
\end{itemize}





\paragraph{5.e}

\begin{itemize}
	\item ¿Es posible definir una distinta para cada usuario?
	
	\item Si, cada usuario tiene definida una shell de inicio que se guarda en el archivo \texttt{/etc/passwd}.
\end{itemize}



\paragraph{6.a}

\begin{itemize}
	\item ¿Qué es el sistema de archivos en Linux?
	
	\item .
\end{itemize}



\end{document}